\chapter{仓库像分房：flag、env、Content ID}
\label{ch:rooms}

\section{Content ID 与刀口是两件事}

\begin{metaphor}[香肠切片 vs 称重贴纸]
刀口决定``切成哪几段''；Content ID（SHA256）是给每段贴的\textbf{指纹贴纸}。\\
换刀法 $\Rightarrow$ 段的起止变了 $\Rightarrow$ 贴纸集合几乎重来。\\
哈希算法本身没换，但\textbf{被哈希的明文区间}变了。
\end{metaphor}

\begin{figure}[htbp]
\centering
\begin{tikzpicture}[font=\small]
  \draw[thick] (0,1.2) rectangle (10,1.8);
  \foreach \x in {2.5,5.5,8} {\draw[CutRed,very thick] (\x,1.15)--(\x,1.85);}
  \node[above] at (5,1.85) {刀口（CDC）};
  \foreach \x/\lab in {1.2/H1, 4/H2, 6.7/H3, 9/H4} {
    \node[softbox=OkGreen, minimum width=1.3cm, minimum height=0.55cm] at (\x,0.3) {\scriptsize\lab};
  }
  \node[font=\footnotesize, gray] at (5,-0.35) {每段明文 $\to$ 一张 SHA256 贴纸};
\end{tikzpicture}
\caption{同一文件、不同刀法 $=$ 不同贴纸集合。}
\label{fig:cid-vs-cut}
\end{figure}

\section{四个独立房间（feature flag）}

\begin{figure}[htbp]
\centering
\begin{tikzpicture}[font=\small]
  \foreach \i/\name/\flag/\col in {
    0/Hom/0x20000/Gen3Color,
    1/Chain/0x40000/Gen4Color,
    2/TopoPH/0x80000/Gen5Color,
    3/Native/0x100000/Gen51Color} {
    \node[softbox=\col, minimum width=2.5cm, minimum height=1.5cm]
      (r\i) at (\i*3.0,0)
      {\textbf{\name}\\\texttt{\flag}\\独立钥匙};
  }
\end{tikzpicture}
\caption{每个 flag 是一把房门钥匙；不能把家具抬去隔壁继续用。}
\label{fig:four-flags}
\end{figure}

\section{进门口令（环境变量）}

\begin{center}
\begin{tabular}{@{}llp{5.5cm}@{}}
\toprule
口令 & 值 & 人话 \\
\midrule
\texttt{EBBACKUP\_CDC\_ALGO} & topocdc/... & 今天进哪间房 \\
\texttt{TOPOPH\_KERNEL} & tri/ph & TopoPH 用翻面还是连通滤网 \\
\texttt{TOPOPHN\_KERNEL} & tri/ph & Native 用 Tri 还是 PH \\
\texttt{TOPOCHAIN\_PARALLEL\_SCAN} & 1 & Chain 允许多人同时找刀 \\
\bottomrule
\end{tabular}
\end{center}

\begin{pitfall}[互斥]
同一进程不能``同时住四间房''。口令互相打架会在开仓/备份时直接 Fail。
\end{pitfall}

\section{variant：房间里的床型}

同房可能有小差异（床型）：
\begin{itemize}[nosep]
  \item Hom 房：variant 0=标准珠串；1=旧 Tri（废弃展品）；
  \item Chain 房：variant 2；
  \item TopoPH：3=翻面滤网，4=连通滤网；
  \item Native：5=Tri-N，6=PH-N。
\end{itemize}

\section{真名对照}

\texttt{kBackupFeatureTopo*}；\texttt{topo\_variant}；\texttt{CdcTopo*Enabled()}；Pipeline 按 Native$\to$PH$\to$Chain$\to$Hom 优先问路。
